\section{Introduction}

Repeated probability surveys describe a population by a distribution rather
than a mean. For an ordinal item the cumulative distribution function records
the share of the population below each response threshold, and a survey
programme covering many countries or many regions produces one such curve per
population per round. These curves are increasingly used as inputs to a second
analysis: information from the observed populations is used to construct a
prediction band for a further one.

The inputs to that second analysis are estimates. Each carries sampling error
from its own complex design, and a band calibrated on the discrepancies between
those estimates and a prediction centre inherits that error. This paper asks
what such a band guarantees, about which object, and under which assumption.

\paragraph{Two targets.} Write $\Ftilde_c$ for population $c$'s survey estimate
and $F_c$ for the population distribution behind it. A band may be asked to
cover a further \emph{survey estimate} $\Ftilde_{K+1}$, which we call target
T1, or the \emph{population distribution} $F_{K+1}$, target T2. The two are not
interchangeable. Exchangeability of the observed curves delivers T1 in finite
samples by the ordinary conformal rank argument. It delivers nothing about T2
on its own, because the calibration scores contain sampling error that the
latent target does not.

Reaching T2 requires additional information, and there are two routes. A bound
on the target's own sampling error gives a shape-free enlargement of the T1
band. Alternatively, subtracting an estimate of sampling variance from the
calibration spread gives a narrower band, at the price of a distributional
assumption: that the standardised observed and latent curves follow a common
law. Correct coordinate variances do not imply it.

\paragraph{What this paper adds.} We do not claim novelty for the conformal
rank argument, for the observation that measurement error is not identified
from a convolution, or for simultaneous scoring of a curve. The contribution
is to connect these to survey-estimated distribution curves and to characterise
when the variance-subtraction route is reachable in practice.

The characterisation turns on a coupling that is not new as algebra but that
appears not to have been drawn out for a calibration procedure. An adaptive
implementation of this correction gates on two quantities: whether there is
enough design noise to be worth removing, and whether the remaining scale is
determined precisely enough to trust. These read as independent requirements,
and a feasibility region drawn from them looks like a corner of a plane. They
are not independent. The deconvolution shrinks the
target scale in proportion to the design share, and the reliability diagnostic
is inflated by exactly that shrinkage. Across our survey applications the
correlation between the design share and the diagnostic's inflation above its
sample-size floor is \RegCorrRhoRatio. \textbf{Design noise large enough to be
worth removing is, by that very magnitude, large enough to leave the remainder
poorly determined.} The feasible region is a narrow diagonal band, not a corner.

This displaces the account that accompanies the implementation. Its stated
obstacle is a population-count floor of ninety-four, derived from the
diagnostic's sampling term. In our regional applications that floor is cleared in
\RegFloorCleared{} of \RegConfigs{} configurations, at population counts up to
\RegKmax, and the correction still activates in only \RegActivations. The
floor is not what binds.

\paragraph{Empirical characterisation.} At the national unit of the European
Social Survey the design share is \NatRhoLo--\NatRhoHi{} on full samples and
reaches \NatScanHi{} in a subgroup scan; the need gate opens in none of
\NatConfigs{} configurations. The correction is unnecessary there. One unit
down, on regions within the same survey, the design share runs
\RegRhoLo--\RegRhoHi{} and the correction becomes both needed and, in a few
configurations, feasible. Where it activates it narrows the band by
\WidGainLo--\WidGainHi{} percent against the uncorrected alternative, so the
boundary tells a practitioner where to look rather than telling them to stop. Which configurations depends on the variance
estimator: a plain $m$-of-$m$ resample opens the reliability gate in
\VerMofmOpen{} of \VerMofmN{} arms and a Rao--Wu--Yue rescaled resample in
\VerRwyOpen{} of \VerRwyN, with the frozen threshold falling inside that gap.

Refining the calibration unit is the natural response to a population-count
requirement, and it works only in one direction. Cutting the survey into regions yields \DomRegK{} populations against a
requirement of \DomRegReq, within reach; cutting the same survey by sex yields
\DomSexK{} against \DomSexReq, because men and women have nearly the same trust
distribution and the split adds sampling noise without adding dispersion. Across
\DomConfigs{} configurations spanning \DomTypes{} domain definitions the
boundary predicts the gate in \DomCorrect{} of them, and \DomOpen{} open.

Finally, a simulation isolates the distributional assumption from
scale estimation. With exact coordinate variances throughout, mismatched
correlation structures cost simultaneous coverage in a direction that is
predictable for cumulative distribution functions, and the cost grows with the
number of coordinates carried simultaneously rather than shrinking with the
number of calibration populations.
